\chapter{Conclusion}

This thesis investigated whether interacting quantum-dot chains can be numerically optimized to realize Majorana-like low-energy states and whether the resulting configurations support a useful braiding protocol. The optimization reproduces the known non-interacting two-dot sweet spot and, for longer chains, identifies more inhomogeneous parameter profiles with strong central structure in the onsite energies and pairing amplitudes. In particular, the recurring interior detuning is correlated with favorable edge-localization diagnostics and suppressed bulk Majorana polarization, although a controlled parameter scan would be required to establish whether it is directly responsible for this behavior.

Among the tested interacting configurations, the weakly interacting three-dot system provides the strongest overall braiding candidate. It combines a small ground-state splitting, a large excitation gap, edge-localized Majorana polarization, and strongly suppressed bulk polarization. Stronger interactions and larger systems were generally more difficult to optimize within the searches performed here. However, the optimization combines local gradient-based updates, repeated restarts, and basin hopping, and the finite numerical search cannot guarantee that the global minimum was found. The optimized configurations should therefore be interpreted as promising parameter regimes rather than proof of perfectly robust Majorana states.

Within the lowest projected manifold, the optimized three-dot system reproduces the expected Majorana exchange with high accuracy. The projected microscopic junction operators are dominated by the intended Majorana bilinears, and the energy-level Majorana reference model, projected microscopic junction, and projected local-operator models give nearly identical braid diagnostics. This demonstrates that the optimized interacting PMM system can support the expected exchange behavior within the controlled low-energy projected model.

The sector-resolved analysis shows that the self-target braid errors remain low across all identified sectors for the non-interacting and weakly interacting configurations. The selected stronger-interaction local-operator constructions contain larger and more sector-dependent errors. However, the stable $U=2.0$ energy-level baseline and irregular $U=1.0$ result show that interaction strength alone does not explain the observed ordering. The quality of the selected optimized configuration and the distinction between the energy-level and local exchange constructions are also important.

The block-matching analysis shows that the projected local operators can be represented essentially exactly within the energy-level Majorana basis. However, different internal blocks can follow opposite exchange orientations relative to the fixed energy-level convention. A higher-dimensional sector therefore does not necessarily implement one uniform logical gate, and predictable braiding across such a sector requires knowledge or control of the occupied block.

Overall, optimization can identify interacting microscopic PMM configurations that reproduce effective Majorana exchange within controlled projected models. Extending this behavior beyond one known low-energy block requires additional control over the occupied block and exchange orientation, as well as a common-target comparison to determine how accurately projected local channels reproduce the energy-level exchange. Natural next steps are to improve the optimization coverage, test the role of controlled interior detuning, and include disorder, noise, finite-temperature effects, and nonadiabatic driving in the braiding analysis.
